\documentclass[a4paper,12pt]{ctexart}
\usepackage{xcolor,graphicx,geometry,array,hyperref,booktabs}
\hypersetup{colorlinks=true,linkcolor=blue!70!black}
\geometry{top=1.8cm,bottom=1.8cm,left=1.8cm,right=1.8cm}
\linespread{1.35}
\newenvironment{thinkbox}{\vspace{0.3cm}\begin{center}\begin{minipage}{0.88\textwidth}\colorbox{blue!5}{\begin{minipage}{0.94\textwidth}\vspace{0.15cm}\textbf{【想一想】}}\vspace{0.15cm}\end{minipage}\end{minipage}\end{center}\vspace{0.3cm}}{}
\newenvironment{funfact}{\vspace{0.2cm}\begin{center}\begin{minipage}{0.88\textwidth}\fbox{\begin{minipage}{0.92\textwidth}\vspace{0.1cm}\textbf{【有趣的小知识】}}\vspace{0.1cm}\end{minipage}\end{minipage}\end{center}\vspace{0.2cm}}{}
\newenvironment{figplaceholder}[1]{\vspace{0.2cm}\begin{center}\fbox{\begin{minipage}{0.82\textwidth}\centering\textbf{【插图位置】}\\[0.2cm]#1\end{minipage}}\end{center}\vspace{0.2cm}}{}

\title{\Huge\textbf{航空航天小故事}}
\date{}
\begin{document}\maketitle\thispagestyle{empty}

\section{火箭是怎么飞上天的？——牛顿第三定律的极致应用}

"火箭是怎么飞上天的？"小明问。

"因为它在向下喷东西。"爸爸答。

"向下喷东西就能向上飞？"小明不信。

"你吹过气球吗——你往气球里吹满气，捏住口，然后松手——气球会嗖地一下飞出去，乱撞。为什么？因为气球里的气体从口子喷出去，把气球推向了相反的方向。火箭的原理完全一样——只是它喷的是高温燃气，喷得极快（每秒几公里），而且喷的口子——叫喷管——是经过精密设计的，让气体的速度最大化。"

这就是\textbf{牛顿第三定律}——作用力等于反作用力（大小相等、方向相反）。但这个定律的应用有一个关键点——不需要外部介质。火箭在太空中（没有空气）也能飞——因为它不需要"推空气"，它只需要"推自己喷出去的那团燃气"。

\begin{figplaceholder}
插图：火箭的喷气原理——燃烧室→喷管→高速燃气喷出，同时火箭向反方向推进
\end{figplaceholder}

\textbf{多级火箭：}为什么火箭要做成"好几节"（级）？因为越往上飞重量越重要。一级火箭烧完燃料后——把空壳丢掉（越级分离）——剩下的部分更轻了——更容易加速。长征五号（"胖五"）是两级火箭，加上助推器——一级分离后，只剩下芯级继续工作。

\textbf{中国的火箭和天宫空间站：}2022年中国空间站"天宫"全面建成。长征五号B负责发射核心舱和实验舱——它的近地轨道运载能力约25吨（是长征二号F的3倍以上）。

\begin{thinkbox}
如果你要发射一颗卫星到太空——你希望火箭越快越好还是越慢越好？——越快越好，因为地球引力一直在拉卫星下来。火箭的速度必须达到约7.9km/s（第一宇宙速度）才能进入近地轨道。如果火箭在半空中停了——卫星会坠回地球。
\end{thinkbox}

\begin{funfact}
2024年SpaceX的星舰（Starship）进行了第三次轨道试飞——这是人类历史上最大的火箭（高121米，比自由女神像还高）。它的第一级（助推器）由33台发动机同时点火——产生的推力相当于一架波音747客机发动机的"8万倍"。中国也在研发可回收重复使用火箭——长征九号深度变推液氧甲烷发动机（200吨级）正在研制中。
\end{funfact}

\newpage

\section{宇航员在太空怎么生活？——失重不是没有引力}

天宫空间站里的宇航员飘来飘去——东西也飘来飘去。小明问："太空中没有地球引力吗？"

"不是的。地球引力在空间站那个高度（约400公里）仍然非常强——大约是地球表面的90\%。如果引力消失了，空间站不会飘在空中——它会沿直线飞走。宇航员失重不是因为引力消失了——而是因为他们和空间站一起在做'自由落体'。"

\textbf{什么是"自由落体"？}

空间站以约7.9km/s的速度绕地球飞行——这个速度让它总是在"掉向地球"，但因为它同时在水平方向快速移动，所以它"掉"的弧度和地球的曲率一致——永远"掉"不到地上。空间站里的宇航员和空间站一同"掉落"——所以他们感觉不到重量。

\begin{figplaceholder}
插图：空间站的轨道——向心力（地球引力）和离心力的平衡——标注"自由落体但永远掉不到地上"
\end{figplaceholder}

\textbf{在太空怎么吃饭、睡觉？}
\begin{itemize}
  \item 吃饭：食物通常是冻干或者罐头——加水就能吃。面包不行——面包屑会在舱内乱飘。所以宇航员吃的是墨西哥卷饼（没有面包屑）
  \item 睡觉：宇航员可以"飘"着睡——没有一个方向是"上"。但如果飘着睡，醒来可能发现自己撞到了天花板。所以他们通常把自己装进一个睡袋、固定在墙上
  \item 上厕所：最复杂的生活任务——没有重力，水和尿液不会"往下流"。太空马桶用气流把排泄物吸到容器里。尿液被回收处理成饮用水（是的——宇航员喝自己过滤后的尿液）
  \item 锻炼：失重状态下肌肉不用"对抗重力"——时间长了会萎缩。宇航员每天必须锻炼2小时（跑步机/自行车/举重——所有器械都用特殊方式固定才能用）
\end{itemize}

\begin{thinkbox}
如果在太空中失重，你的身高会发生什么变化？——会变高（约3-5厘米）。因为脊柱在没有重力压迫的情况下"伸展"了。返回地球后几周内恢复正常。这就是宇航员"长高"的秘密。
\end{thinkbox}

\begin{funfact}
国际空间站上的宇航员一天能看16次日出和日落——因为空间站每90分钟绕地球一圈。从太空看地球——没有国境线、没有边界——只有一块蓝色的球体。很多宇航员说，从太空回望地球后，他们对"人类是一个整体"这件事有了完全不同的感受——这就是宇航员的"总观效应"（Overview Effect）。
\end{funfact}

\newpage

\section{飞机为什么能飞起来？——伯努利原理与反作用力}

在intro故事第二课里，小明问了："飞机那么重——为什么能飞在天上不掉下来？"

\textbf{两种解释：}

\textbf{解释一（伯努利原理）：}机翼上表面是弧形的，下表面是平的。空气流过上表面时速度更快——流速快的地方压强低，流速慢的地方压强高。机翼下方压强高于上方——产生向上的升力。

这个解释在教科书里出现最多。但它不是完整的答案——它只能解释一部分升力。

\textbf{解释二（牛顿第三定律）：}飞机机翼以一定的角度（攻击角）向前运动，把空气向后下方"推"——空气的反作用力就把机翼向前上方推。升力主要来自这个"反作用力"——而不是伯努利原理的压强差。

实际上，伯努利原理和牛顿定律协同工作——二者都对升力有贡献。

\begin{figplaceholder}
插图：飞机机翼的剖面——标注上表面弧形（高速/低压）、下表面平直（低速/高压）、攻击角、气流方向
\end{figplaceholder}

\textbf{纸飞机实验：}不同折法的纸飞机飞行表现完全不同。宽机翼产生更大升力但阻力也大（飞得慢）；窄机翼阻力小但升力也小（飞得远但需要速度）。这就是空气动力学中的"升阻比"——飞机设计师永远在"升力"和"阻力"之间寻找最优平衡。

\begin{thinkbox}
尝试折一架纸飞机——什么形状飞得最远？什么形状可以做"回旋镖"飞回你手里？试一下然后思考——你的设计改变了什么参数？
\end{thinkbox}

\begin{funfact}
波音747的机翼长度约60米——比莱特兄弟第一次飞行的距离（36米）还要长。莱特兄弟的"飞行者一号"在1903年12月17日飞了12秒、36米——到2024年，全球每天约10万个航班、运送约1300万名乘客。飞机已经成为现代人类最安全的交通工具——按每10亿公里计算的死亡率，飞机是汽车的1/200。
\end{funfact}

\end{document}
